\documentclass[11pt,a4paper]{article}
\usepackage[margin=25mm]{geometry}
\usepackage{amsmath,amssymb,amsthm,hyperref}
\setlength{\emergencystretch}{2em}
\newtheorem{lemma}{Lemma}
\newtheorem{theorem}[lemma]{Theorem}
\newcommand{\gen}[1]{\langle#1\rangle}
\title{Transitive label budgets on six residual fixed points}
\author{Human Compression project: L044 continuation}
\date{}
\begin{document}\maketitle
Use the critical singleton-end setting of the accepted residual-deficit
proof: $G=\gen{\rho_0,\ldots,\rho_{r-1}}\leq A_n$ is transitive,
is a connected string C-group, and has every maximal omission intransitive;
$G_0$ has orbits $\{a\}$ and its complement;
$\rho_0=(a\ b)\delta$, $K=G_{>1}$, $R=G_{>2}$ and $\sigma=\rho_1$.
The odd involution $\delta$ has $t$ pairs, with block set $\mathcal T$,
union $P$, and an even
fixed set $F$ of size $f$ in $\Omega\setminus\{a,b\}$.
Assume $n=2r\geq14$, so $K$ has $d=t+f/2-1$ labels.
All block-image omission-orbit statements below use the complete
accepted actual-deficit proof. Whiston's bound and the complete
odd-degree maximal-string theorem, including degree five, remain paid.

\begin{lemma}[Single-transposition restrictions]\label{single}
If $K$ is transitive on $F$, $f\geq4$, and all nonidentity distinguished
restrictions to $F$ are single transpositions, they are exactly the
$f-1$ successive edges of a natural path, each with one label.
\end{lemma}
\begin{proof}
Their edge graph is connected. Three distinct edges incident to one
vertex would have pairwise noncommuting labels, impossible in a string.
A cycle is likewise impossible: labels chosen on its distinct edges
would require its line graph, a cycle, to embed in the path of label
positions. Hence the simple edge graph is a path.

Choose a label on each edge of that path. Their positions form a path
in the integer line, so are consecutive in one of the two orders.
Any repeated label on an interior edge would have to be adjacent to
both chosen neighboring labels, whose unique common neighbor is the
already occupied position. For an end edge, its neighboring edge is
interior because $f\geq4$; both positions adjacent to that neighboring
label are already occupied. A second end-edge label is also impossible.
\end{proof}

\begin{lemma}[Minimum transitive fixed-set budget]\label{budget}
Let $E$ be a set of $h$ distinguished $K$ labels whose restrictions
are transitive on $F$. Then $h\geq f/2$.
If equality holds and $t\geq3$ is odd, the other $t-1$ pair images
generate the ordinary $S_t$ path, at consecutive positions $J$,
and every label of $E$ has trivial pair image.
\end{lemma}
\begin{proof}
Deleting a label outside $E$ leaves the fixed-set action transitive.
Its pair image cannot be generated by the other pair images, as that
would leave the full block-orbit partition unchanged. In particular
the $d-h$ pair images outside $E$ are independent. Whiston gives
$d-h\leq t-1$, or $h\geq f/2$.

At equality those $t-1$ images generate $S_t$ and form its ordinary
path by the accepted odd-degree maximal-string theorem. For $t=3$
this just says that two independent involutions in $S_3$ are distinct
transpositions. Consecutive noncommutations force consecutive positions
in the original string. Omitting any path mark leaves two pair orbits.
If an $E$ label joined them, the full omission would be transitive both
on pairs and on $F$, contradicting the actual orbit deficit. Thus every
$E$ label preserves every path cut; their common stabilizer fixes every
pair-block. This proves the last assertion.
\end{proof}

\begin{lemma}[The first label cannot be an extra]\label{near}
Let $E$ consist of $h=f/2\geq2$ distinguished labels, transitive on $F$,
and suppose $t\geq3$ is odd.
Then $\rho_2\notin E$.
\end{lemma}
\begin{proof}
Suppose $\rho_2\in E$. Then $R$ contains all of $J$, so is transitive
on $\mathcal T$ and fixes no point of $P$. Because $K$ fixes $b$,
transitivity of $G_0=\gen{\sigma,K}$ forces $\sigma(b)\ne b$.
Also $\sigma(b)\ne a$, since $\sigma$ fixes $a$. Centralization
shows that $z=\sigma(b)$ is fixed by $R$; hence $z\in F$.

Partition the positions of $E$ into maximal intervals. Distinct
intervals commute, by the string relations. Any interval not containing
$\rho_2$ is contained in $R$, and its subgroup fixes $z$. It also
commutes with the subgroup on all the other intervals. Consequently
it fixes the entire $\gen{E}$-orbit of $z$: each orbit point can be
written $z^u$ with $u$ in those other intervals, because the interval
itself fixes $z$. This orbit is $F$, by transitivity. Removing even one
label of that interval would then leave the same transitive action on
$F$, contradicting the minimum $h$-label budget.

There is therefore only one $E$ interval, beginning at $\rho_2$.
All $J$ labels follow it. Since $|J|=t-1\geq2$, its last label
$\lambda$ is at distance at least two from every $E$ label, and
centralizes $\gen{E}$. It belongs to $R$, so fixes $z$, and hence
fixes all of $F$ pointwise. It also fixes $a,b$.

On the other hand, $\lambda$ acts nontrivially on every $R$-orbit
$O\subseteq P$. Such an orbit meets every pair-block, since $R$ is
transitive on $\mathcal T$; the nontrivial pair transposition of
$\lambda$ therefore moves points of $O$. The permutation $\sigma$
transports $R$-orbits equivariantly. No such $O$ can be sent outside
$P$, where $\lambda$ fixes every point. Thus $\sigma$ preserves $P$.
So does $K$, contradicting transitivity on the complementary $G_0$
orbit. This proves the lemma.
\end{proof}

\begin{theorem}[The strict transitive label budget]\label{strict}
Suppose $f\geq4$ and $t\geq3$. Every transitive fixed-set subfamily
has at least $f/2+1$ labels. In particular no three distinguished
labels can be transitive on a six-point fixed set, including
$t=3$, $n=14$.
\end{theorem}
\begin{proof}
Lemma~\ref{budget} gives $h\geq f/2$; suppose equality holds.
By Lemma~\ref{near}, $\rho_2$ belongs to $J$. Its consecutive path
therefore occupies the initial positions, and all $E$ labels lie in
$R$. Thus $F$ is an $R$-orbit. If $\sigma$ preserved it, then $G_0$
would preserve it too. Hence $\sigma(F)=O\subseteq P$; the singleton
orbits $\{a\},\{b\}$ cannot be its image.

Equivariance makes $\gen{E}$ transitive on $O$. But every $E$ label
has trivial pair image: each individual pair-block is invariant under
$\gen{E}$. Its orbits in $P$ consequently have size at most two,
contradicting $|O|=|F|=f\geq4$. Equality is impossible.
\end{proof}

\paragraph{Exact six-point consequence.}
A distinguished triple transposition on a transitive six-point $F$
can be extended from its three disjoint edges to a spanning tree using
two further edges. At most three labels are then transitive on $F$,
contrary to Theorem~\ref{strict} and the minimum budget. Thus no triple
transposition occurs. If every nonidentity restriction is a single
transposition, the separate path-label lemma and natural fixed-path
exclusion apply. The remaining transitive six-point case has at least
one double transposition and minimum transitive label number exactly
four: a spanning tree extending its two edges uses at most three
further labels, while the lower bound and Theorem~\ref{strict} exclude
three. This last case is not resolved here.

\paragraph{Attribution and boundary.}
The source is Cameron--Fernandes--Leemans--Mixer,
\emph{Highest rank of a polytope for $A_n$}, arXiv:1605.09173v1,
physical pages38--44. The path-label observation and the accepted four-fixed argument of collaborator
\texttt{rewrite\_termination} supply the underlying equivariant-orbit
method; the full preceding actual-deficit argument remains paid.
Lemmas~\ref{single}--\ref{budget} are unchanged from the preceding
label-budget edition; its weaker three-label theorem is strengthened here. The fixed point of $R$, interval-minimality and distant
path-label obstruction are the present continuation. The final
individual-pair orbit observation was also explicitly supplied by the
root collaborator while this successor was being completed. No numerical
search, projected IP, smaller-degree alternating induction or new
formal theorem is used. The intransitive six-point actions, four-label
case, larger fixed sets and other principal L044 branches remain open
in this recovery. No whole-original or question credit is claimed.
\end{document}
